\section{A preserved constraint on posterior laws}\label{sec:cone}

The second moment $\E A(p)$ measures information about the input, but the
estimate we need for products also involves one higher-order statistic.
We now give this statistic explicitly. Its expectation will remain
nonnegative throughout the recursion, allowing a strict improvement on
second-moment subadditivity.

For $x=Lp$, with $L$ as in Section~\ref{sec:experiments}, set
\[
 B(p)=x_1x_2x_3,\qquad
 C(p)=x_1^2x_2^2+x_1^2x_3^2+x_2^2x_3^2.
\]
Together with $A(p)=x_1^2+x_2^2+x_3^2$, these polynomials are invariant under
permuting the four entries of $p$. Indeed, such a permutation acts on $x$
by a permutation of its coordinates and an even number of sign changes.
Define $H,F,G$ by the following coefficient table, whose entries are divided
by $1000$:
\begin{equation}\label{eq:polynomial-table}
\begin{array}{c|rrrrrrrrr}
 &1&A&B&A^2&C&AB&A^3&AC&B^2\\ \hline
 H&0&0&1000&-100&-150&300&-20&70&-510\\
 F&0&4000&-3000&700&-4000&-800&0&0&0\\
 G&1000&-1610&-3100&773&4000&-3980&0&0&0
\end{array}
\end{equation}
Thus, for example, $H=B-A^2/10-3C/20+3AB/10-A^3/50
+7AC/100-51B^2/100$. We also write $H(x)$ when its three coordinate
arguments, rather than the probability vector, are supplied.

For $p,q\in\Deltafour$ and $\pi\in\Sym$, put
\[
 U_i^{\pi}=4p_iq_{\pi(i)},\qquad
 z_\pi=\sum_{i=1}^4 U_i^{\pi},\qquad
 P^{\pi}=U^{\pi}/z_\pi\quad(z_\pi>0).
\]
For a continuous function $\phi$ on $\Deltafour$, use the notation
\begin{equation}\label{eq:permutation-average}
 [z\phi(P)](p,q)=\frac1{24}\sum_{\pi\in\Sym}
                 z_\pi\phi(P^{\pi}),
\end{equation}
with each summand set to zero when $z_\pi=0$. Finally, define
\begin{equation}\label{eq:cubic-saving}
 D_0(p,q)=\frac{A(p)A(q)\bigl(A(p)+A(q)\bigr)}{1000}.
\end{equation}

\begin{proposition}[Polynomial inequalities]\label{prop:polynomial-inequalities}
For every $p,q\in\Deltafour$, the polynomials in
\eqref{eq:polynomial-table} satisfy $F(p)\geq0$, $G(p)\geq0$, and
$H(p)\geq-1/2$. For $x=Lp$ and $0\leq t\leq1$,
\begin{equation}\label{eq:H-scaling}
 H(tx)\geq t^3H(x).
\end{equation}
Moreover,
\begin{align}
 [zA(P)](p,q)
 &\leq A(p)+A(q)-H(p)F(q)-F(p)H(q)-D_0(p,q),
 \label{eq:A-product}\\
 [zH(P)](p,q)
 &\geq H(p)G(q)+G(p)H(q).
 \label{eq:H-product}
\end{align}
\end{proposition}

Sections~\ref{sec:polynomial-certificates} and
\ref{sec:finite-verification} prove this proposition using exact
polynomial sign certificates. We first show how it controls posterior laws;
this identifies the purpose of each inequality before the coefficient
calculation begins.

Let $\mathcal C$ be the class of symmetric posterior laws $\mu$ for which
\begin{equation}\label{eq:cone-condition}
 \E_{\mu} H(p)\geq0.
\end{equation}
The condition is linear in the law. It does not require $H(p)\geq0$ at
every posterior: for example, $p=(1/2,1/2,0,0)$ has $H(p)=-3/25$.
The completely informative experiment and the uninformative experiment
belong to $\mathcal C$, because
\begin{equation}\label{eq:initial-H}
 H(e_i)=\frac{13}{100}\quad(1\leq i\leq4),\qquad H(\unif)=0.
\end{equation}
Here the informative experiment has posterior law uniform on the four
point masses $e_i$.

\begin{proposition}[Closure and moment saving]\label{prop:cone-closure}
The class $\mathcal C$ is preserved by a channel step with
$\lambda\in[0,1]$, by conditionally independent products of experiments,
and by mixtures whose mixing label is observed and independent of the
input. If $\mu,\nu\in\mathcal C$ are the marginal posterior laws of two
experiments, then their product posterior law $\mu\star\nu$ satisfies
\begin{align}
 \E_{\mu\star\nu} A
 \leq{}& \E_{\mu}A+\E_{\nu}A\notag\\
 &-\frac{\E_{\mu}A^2\,\E_{\nu}A+
                \E_{\mu}A\,\E_{\nu}A^2}{1000}.
 \label{eq:moment-saving}
\end{align}
In particular, when $\E_{\mu}A=\E_{\nu}A=b$, the subtracted term is at
least $2b^3/1000$.
\end{proposition}

\begin{proof}
By Lemma~\ref{lem:channel}, a channel step scales $x=Lp$ by $\lambda$
without changing the marginal law of the observation. Consequently
\eqref{eq:H-scaling} gives
\[
 \E H(\lambda x)\geq\lambda^3\E H(x)\geq0.
\]
The resulting law remains symmetric.

For a product, Lemma~\ref{lem:bayes-product} expresses its posterior
expectations as $\E_{\mu\otimes\nu}[z\phi(P)]$. The expectation here is
over independent marginal posterior laws, before multiplication by the
density $z$. Since $\nu$ is symmetric, it is unchanged by permuting its
four coordinates. We may therefore replace $z\phi(P)$ by its permutation
average whenever $\phi$ is $A$ or $H$. Integrating
\eqref{eq:H-product} and using independence gives
\[
 \E_{\mu\star\nu}H
 \geq \E_{\mu}H\,\E_{\nu}G+
        \E_{\mu}G\,\E_{\nu}H\geq0.
\]
Here $G\geq0$ pointwise, whereas the two expectations of $H$ are
nonnegative by membership in $\mathcal C$. Symmetry is preserved by
Lemma~\ref{lem:bayes-product}.

Similarly, integration of \eqref{eq:A-product} yields
\[
 \E_{\mu\star\nu}A
 \leq \E_{\mu}A+\E_{\nu}A
       -\E_{\mu}H\,\E_{\nu}F
       -\E_{\mu}F\,\E_{\nu}H
       -\E_{\mu\otimes\nu}D_0(p,q).
\]
The two terms involving $F$ can be discarded because $F\geq0$ and
$\mu,\nu\in\mathcal C$. Expanding \eqref{eq:cubic-saving} gives
\eqref{eq:moment-saving}. Since $A\geq0$, Jensen's inequality gives
$\E_{\mu}A^2\geq(\E_{\mu}A)^2$ and the analogous bound for $\nu$,
which proves the last assertion.

Finally, conditioning on an observed mixing label independent of the input
leaves the uniform prior unchanged. The posterior law of the mixed
experiment is the corresponding mixture of posterior laws. Symmetry and
\eqref{eq:cone-condition} are both preserved by this mixture.
\end{proof}
